%% ma: L12-B14-0142
%% lop: 12
%% chuong: 5
%% bai: 14
%% dang: 12.14.11
%% loai: TLN
%% muc: TH
%% dap_an: "-1"
%% nguon: "Ảnh giáo viên gửi 25/11/2026 (tài liệu ôn tập Phương trình mặt phẳng), Phần III Câu 1"
%% kien_thuc: "Bài 14 (mặt phẳng song song với mặt phẳng cho trước)"
%% trang_thai: chinh-thuc
\begin{ex}
Trong không gian $Oxyz$, biết mặt phẳng đi qua điểm $M(1;0;0)$ và song song với mặt phẳng $(\alpha)\colon 2x-y-1=0$ có dạng $-2x+ay+bz-c=0$, tính $a+b+c$.
\shortans{-1}
\loigiai{Mặt phẳng song song $(\alpha)$ có dạng $2x-y+d=0$ ($d\ne-1$). Qua $M$: $2+d=0\Rightarrow d=-2$, nên mặt phẳng là $2x-y-2=0\Leftrightarrow-2x+y+2=0$. Đồng nhất với $-2x+ay+bz-c=0$: $a=1$, $b=0$, $c=-2$. Vậy $a+b+c=-1$.}
\end{ex}
